\chapter{Evaluation}
\label{ch:evaluation}

This chapter evaluates your contribution against the research questions stated in Chapter~\ref{ch:introduction}. Define clear evaluation criteria before presenting results — what does success look like?\footnote{A good evaluation is systematic: it states what is being measured, how, and why that measure is meaningful.}

%%%%%%%%%%%%%%%%%%%%%%%%%%%%%%%%%%%%%%%%%%%%%%%%%%%%%%%%%%%%%%%%%%%%%%%%
\section{Experimental Setup}
\label{sec:experimental-setup}

Describe the conditions under which your evaluation was conducted. Include datasets, baselines, metrics, and any relevant hardware or software configuration.\footnote{Be precise enough that someone else could replicate your evaluation.}

%%%%%%%%%%%%%%%%%%%%%%%%%%%%%%%%%%%%%%%%%%%%%%%%%%%%%%%%%%%%%%%%%%%%%%%%
\section{Metrics}
\label{sec:metrics}

Define the metrics you use to assess your approach. Explain why each metric is appropriate for your research questions.

%%%%%%%%%%%%%%%%%%%%%%%%%%%%%%%%%%%%%%%%%%%%%%%%%%%%%%%%%%%%%%%%%%%%%%%%
\section{Results}
\label{sec:eval-results}

Present the evaluation results. Use tables and figures to summarise comparisons across conditions or baselines.

\begin{table}[ht]
  \centering
  \begin{tabular}{lcc}
    \toprule
    \textbf{Method} & \textbf{Metric A} & \textbf{Metric B} \\
    \midrule
    Baseline   & 0.00 & 0.00 \\
    Your method & 0.00 & 0.00 \\
    \bottomrule
  \end{tabular}
  \caption{Comparison of your method against the baseline across evaluation metrics.}
  \label{tab:eval-results}
\end{table}

%%%%%%%%%%%%%%%%%%%%%%%%%%%%%%%%%%%%%%%%%%%%%%%%%%%%%%%%%%%%%%%%%%%%%%%%
\section{Discussion}
\label{sec:eval-discussion}

Interpret the evaluation results. Do they support your hypotheses? Where does your approach perform well and where does it struggle? Relate findings back to the research questions.\footnote{If results are mixed or negative, discuss why — this is valuable scientific contribution in itself.}
